\chapter{泛型}
\label{ch:14}

第 13 章的接口，解决的是「一个函数吃多种类型」——但它靠的是运行期的胖引用和类型号，每调一次都有分派开销。

这一章的泛型解决同一个问题，思路完全不同：在编译期为每一种用到的类型，分别生成一份专门的代码。跑起来没有分派，因为它压根就是一开头就定死的。

先看它要治什么病。

\section{问题：换一种类型，逻辑就得抄一遍}

写一个「谁大取谁」的函数。整数版：

\par
\medskip
\begin{lstlisting}[style=magel]
int64 maxOf(int64 a, int64 b)
{
    if (a > b)
    {
        return a;
    }
    return b;
}
\end{lstlisting}
\par
\medskip

字符串也想要一个——逻辑一字不差，只有类型不同：

\par
\medskip
\begin{lstlisting}[style=magel]
String maxOf(String a, String b)
{
    if (a > b)
    {
        return a;
    }
    return b;
}
\end{lstlisting}
\par
\medskip

\texttt{int64}、\texttt{float64}、\texttt{String}……每加一种类型就得抄一份。用第 4 章的重载能跑，但你心里清楚：这三份代码长得一模一样，只是类型不同。 类型之间应该有办法复用。

\section{泛型函数：把「类型」也变成参数}

泛型的写法很简单：在函数名的尖括号里声明一个类型参数 \texttt{T}，然后在原本写类型的地方写 \texttt{T}。

\par
\medskip
\begin{lstlisting}[style=magel]
T maxOf<T>(T a, T b)
{
    if (a > b)
    {
        return a;
    }
    return b;
}

int64 main()
{
    println("{}", maxOf(3, 9));
    println("{}", maxOf("apple", "pear"));
    return 0;
}
\end{lstlisting}
\par
\medskip

\par
\medskip
\begin{lstlisting}[style=out]
9
pear
\end{lstlisting}
\par
\medskip

\texttt{T} 是一个占位符，代表「先不定，等调用时再说」的类型：

\begin{itemize}

\item \texttt{max\allowbreak{}Of(\allowbreak{}3,\allowbreak{} 9)\allowbreak{}} —— 编译器一看实参是 \texttt{int64}，就把 \texttt{T} 当成 \texttt{int64}；

\item \texttt{max\allowbreak{}Of(\allowbreak{}"apple",\allowbreak{} "pear")\allowbreak{}} —— 这次 \texttt{T} 是 \texttt{String}。

\end{itemize}

一个定义，两种类型。 加第三种也不用改 \texttt{max\allowbreak{}Of} 一个字。

\begin{tip}{提示}
调用时 \texttt{T} 是怎么定下来的？靠实参推导——编译器看实参是什么类型，就把 \texttt{T} 绑成什么类型。
\end{tip}

\section{显式写出来，或者让编译器推}

推导不是每次都行。你可以把类型实参显式写出来，写法是把类型放在函数名后的尖括号里：

\par
\medskip
\begin{lstlisting}[style=magel]
T maxOf<T>(T a, T b)
{
    if (a > b)
    {
        return a;
    }
    return b;
}

int64 main()
{
    println("{}", maxOf<int64>(3, 9));
    println("{}", maxOf<String>("apple", "pear"));
    return 0;
}
\end{lstlisting}
\par
\medskip

\par
\medskip
\begin{lstlisting}[style=out]
9
pear
\end{lstlisting}
\par
\medskip

什么时候必须显式？实参推不出唯一答案的时候。 比如 \texttt{int64} 和 \texttt{float64} 混着喂：

\par
\medskip
\begin{lstlisting}[style=magel]
T maxOf<T>(T a, T b)
{
    if (a > b)
    {
        return a;
    }
    return b;
}

int64 main()
{
    println("{}", maxOf(1, 2.0));
    return 0;
}
\end{lstlisting}
\par
\medskip

报错直接点名：

\par
\medskip
\begin{lstlisting}[style=out]
magel：泛型函数 maxOf 的类型实参推导冲突，请显式指定类型实参
note: 第 2 个实参实际为 float64，形参模式为 ...
\end{lstlisting}
\par
\medskip

\texttt{1} 是 \texttt{int64}、\texttt{2.\allowbreak{}0} 是 \texttt{float64}，编译器推不出 \texttt{T} 该是哪一个——Magel 不做隐式转换（第 3 章），所以它宁可停下来让你说清楚，也不自作主张。想让这个调用成立，就写 \texttt{max\allowbreak{}Of<\allowbreak{}float64>\allowbreak{}(\allowbreak{}1.\allowbreak{}0,\allowbreak{} 2.\allowbreak{}0)\allowbreak{}}（注意 \texttt{1} 也要改成 \texttt{1.\allowbreak{}0}，\texttt{float64} 之间才谈得上比较）。

\section{泛型结构体}

结构体也能带类型参数。声明的时候在名字后写 \texttt{<\allowbreak{}>\allowbreak{}}，字段就能用这些参数当类型：

\par
\medskip
\begin{lstlisting}[style=magel]
struct Duo<T, U>
{
    T first;
    U second;
};

int64 main()
{
    Duo<int64, String> d;
    d.first = 1;
    d.second = "a";
    println("{} {}", d.first, d.second);
    return 0;
}
\end{lstlisting}
\par
\medskip

\par
\medskip
\begin{lstlisting}[style=out]
1 a
\end{lstlisting}
\par
\medskip

\texttt{Duo<\allowbreak{}T,\allowbreak{} U>\allowbreak{}} 要两个类型参数。\texttt{Duo<\allowbreak{}int64,\allowbreak{} String>\allowbreak{}} 是一个「第一格是整数、第二格是字符串」的结构。\texttt{Duo<\allowbreak{}String,\allowbreak{} int64>\allowbreak{}} 又是另一个。

\begin{tip}{一个坑}
别给泛型起名 \texttt{Pair}。\texttt{Pair} 是标准库自带的 \texttt{Pair}（随 \texttt{lang} 自动加载，不用 import），你再声明一个同名的泛型，会报「泛型：重复的泛型声明 Pair」——那个报错没提「和标准库撞了」，容易懵。
\end{tip}

泛型结构体的方法也能用 \texttt{T}：

\par
\medskip
\begin{lstlisting}[style=magel]
struct Box<T>
{
    T value;
    T get()
    {
        return value;
    }
    void put(T v)
    {
        value = v;
    }
};

int64 main()
{
    Box<int64> b;
    b.put(41);
    println("{}", b.get());
    Box<String> s;
    s.put("hi");
    println("{}", s.get());
    return 0;
}
\end{lstlisting}
\par
\medskip

\par
\medskip
\begin{lstlisting}[style=out]
41
hi
\end{lstlisting}
\par
\medskip

\texttt{Box<\allowbreak{}int64>\allowbreak{}} 和 \texttt{Box<\allowbreak{}String>\allowbreak{}} 各自有一份 \texttt{put} / \texttt{get}，参数和返回值都变成了具体的类型。

\section{约束：给类型参数提要求}

\texttt{max\allowbreak{}Of<\allowbreak{}T>\allowbreak{}} 里用了 \texttt{a >\allowbreak{} b}。这要求 \texttt{T} 支持比大小——如果你拿一个没有 \texttt{>\allowbreak{}} 的结构体去实例化，它会报错。与其等报错，不如提前声明这个要求。

约束写在类型参数后面，用冒号接一个接口（第 13 章的接口）：

\par
\medskip
\begin{lstlisting}[style=magel]
interface Valued
{
    int64 value();
};

struct Num: Valued
{
    int64 n;
    int64 value()
    {
        return n;
    }
};

int64 read_value<T: Valued>(T item)
{
    return item.value();
}

int64 main()
{
    Num x;
    x.n = 42;
    println("{}", read_value(x));
    return 0;
}
\end{lstlisting}
\par
\medskip

\par
\medskip
\begin{lstlisting}[style=out]
42
\end{lstlisting}
\par
\medskip

\texttt{<\allowbreak{}T:\allowbreak{} Valued>\allowbreak{}} 读作「\texttt{T} 必须是实现了 \texttt{Valued} 的类型」。有了这条，函数体里就能放心地调 \texttt{item.\allowbreak{}value(\allowbreak{})\allowbreak{}}——编译器知道 \texttt{T} 一定有这个方法。多个约束用逗号隔开：\texttt{<\allowbreak{}T:\allowbreak{} A,\allowbreak{} U:\allowbreak{} B>\allowbreak{}}。

约束把「报错时机」提前了。拿一个没实现 \texttt{Valued} 的类型去调，编译器立刻拦下：

\par
\medskip
\begin{lstlisting}[style=magel]
interface Valued
{
    int64 value();
};

int64 read_value<T: Valued>(T item)
{
    return item.value();
}

int64 main()
{
    println("{}", read_value(3));
    return 0;
}
\end{lstlisting}
\par
\medskip

\par
\medskip
\begin{lstlisting}[style=out]
magel：read_value 的类型实参 int64 未实现接口 Valued
\end{lstlisting}
\par
\medskip

一句话说清：约束是写给编译器看的承诺，它保证函数体里那些方法一定存在。

\section{单态化：编译期展开，运行期零开销}

泛型最有意思的地方，是它根本不存在于运行期。

编译器在编译时，为每个实际用到的类型实参，生成一份专门的函数/结构体：\texttt{max\allowbreak{}Of<\allowbreak{}int64>\allowbreak{}} 编译出一份只处理 \texttt{int64} 的代码，\texttt{max\allowbreak{}Of<\allowbreak{}String>\allowbreak{}} 编译出另一份。这个做法叫单态化（monomorphization）——「单一形态化」。

三个直接后果：

\begin{enumerate}

\item 没有运行期分派。 调用 \texttt{max\allowbreak{}Of(\allowbreak{}3,\allowbreak{} 9)\allowbreak{}} 和调用一个手写的 \texttt{max\allowbreak{}Of(\allowbreak{}int64,\allowbreak{} int64)\allowbreak{}} 一样快，机器码是同一份；这跟第 13 章的接口胖引用完全不同（那个要查类型号）。

\item 泛型代码的检查发生在实例化时。 \texttt{max\allowbreak{}Of<\allowbreak{}T>\allowbreak{}} 的函数体里写着 \texttt{a >\allowbreak{} b}，编译器在没看到具体类型之前不会去挑它的毛病；等你拿一个不支持 \texttt{>\allowbreak{}} 的类型来实例化，它才报错——而且报的是那个实例的错：

\end{enumerate}

\par
\medskip
\begin{lstlisting}[style=magel]
T twiceOf<T>(T v)
{
 return v + v;
}

struct P
{
 int64 x;
};

int64 main()
{
 P p;
 p.x = 1;
 println("{}", twiceOf(p).x);
 return 0;
}
\end{lstlisting}
\par
\medskip

\par
\medskip
\begin{lstlisting}[style=out]
magel：没有匹配的运算符接口实现：P + P（由左操作数实现 Add<R,T> 这类左接口，
或由右操作数实现 RightAdd<R,T> 这类右接口）；不会偷偷按数值算
\end{lstlisting}
\par
\medskip

\texttt{twice\allowbreak{}Of<\allowbreak{}int64>\allowbreak{}} 和 \texttt{twice\allowbreak{}Of<\allowbreak{}String>\allowbreak{}} 都没问题（整数 \texttt{+} 整数、字符串 \texttt{+} 字符串），只有 \texttt{twice\allowbreak{}Of<\allowbreak{}P>\allowbreak{}} 这一份不行。

\begin{enumerate}[start=3]

\item 代码会变大。 每个实例都是一份独立的机器码，用一百种类型就生成一百份。这是泛型的代价，换来的是运行期的干净。

\end{enumerate}

\begin{tip}{提示}
顺带一提：\texttt{Pair<\allowbreak{}K,\allowbreak{} V>\allowbreak{}}（14.4 节提过的那个）本身就是泛型结构体——标准库也在用同一套机制。
\end{tip}

\section{默认类型参数}

类型参数可以有默认值，写在参数名后面用 \texttt{=\allowbreak{}} 接：

\par
\medskip
\begin{lstlisting}[style=magel]
struct Wrap<T, U = T>
{
    T a;
    U b;
};
\end{lstlisting}
\par
\medskip

\texttt{Wrap<\allowbreak{}int64>\allowbreak{}} 会写成 \texttt{Wrap<\allowbreak{}int64,\allowbreak{} int64>\allowbreak{}}——\texttt{U} 缺省就跟着 \texttt{T} 走。默认值可以引用前面的参数（\texttt{U =\allowbreak{} T} 合法），但不能引用后面的、也不能是自己。

\begin{tip}{提示}
默认类型参数在初学阶段基本用不到，先把 14.2–14.6 吃透。这一节只是让你在别人的代码里见到 \texttt{struct Foo<\allowbreak{}T,\allowbreak{} U =\allowbreak{} T>\allowbreak{}} 时知道那是什么意思。
\end{tip}

\section{接口还是泛型？}

这两个都能「一个函数吃多种类型」，怎么选？看一遍这张表：

\begin{center}
\begin{tabular}{lll}
\toprule
 & 接口（第 13 章） & 泛型（本章） \\
\midrule
何时定类型 & 运行期 & 编译期 \\
开销 & 胖引用 + 类型号分派 & 无（代码已展开） \\
代价 & 每次调用有分派成本 & 代码体积变大 \\
怎么用 & 类型在运行期才知道（如从容器里取） & 类型在编译期就能定 \\
共同点 & 用接口做约束（\texttt{<\allowbreak{}T:\allowbreak{} Valued>\allowbreak{}}）能结合起来 & 同左 \\
\bottomrule
\end{tabular}
\end{center}

经验法则：

\begin{itemize}

\item 编译期就知道类型 \ensuremath{\rightarrow} 泛型；运行期才定 \ensuremath{\rightarrow} 接口。

\item 两者能合起来用：泛型加约束 \texttt{<\allowbreak{}T:\allowbreak{} Shape>\allowbreak{}} 就是「泛型 + 接口」，既零开销、又能调接口方法。

\end{itemize}

\subsection{编译期整数参数}

泛型支持 \texttt{N:\allowbreak{} int64} 等整数泛型参数，长度参与单态化，不能传入运行期变量。结构体的一维和二维定长数组字段使用这些常量确定布局。

\par
\medskip
\begin{lstlisting}[style=magel]
struct Buffer<T, N: int64>
{
    T data[N];
    int64 len()
    {
        return N;
    }
};
int64 main()
{
    Buffer<int64, 4> buffer;
    buffer.data[0] = 7;
    println("{} {}", buffer.len(), buffer.data[0]);
    return 0;
}
\end{lstlisting}
\par
\medskip

\par
\medskip
\begin{lstlisting}[style=out]
4 7
\end{lstlisting}
\par
\medskip

\texttt{Buffer<\allowbreak{}int64,\allowbreak{}4>\allowbreak{}} 与 \texttt{Buffer<\allowbreak{}int64,\allowbreak{}8>\allowbreak{}} 是不同实例。二维矩阵同理：

\par
\medskip
\begin{lstlisting}[style=magel]
struct Matrix<T, rows: int64, cols: int64>
{
    T data[rows][cols];
};
\end{lstlisting}
\par
\medskip

长度可以有默认值，仍须是可在编译期求值的 int64 常量；运行期长度继续用 \texttt{T a[\allowbreak{}]\allowbreak{} =\allowbreak{} new[\allowbreak{}n]\allowbreak{}}。

\subsection{关联类型}

接口体内的 \texttt{type Item;} 声明关联类型，实现者用 \texttt{type Item =\allowbreak{} T;} 提供具体类型，受约束泛型可以写 \texttt{C.\allowbreak{}Item}。

\par
\medskip
\begin{lstlisting}[style=magel]
interface Container
{
    type Item;
    Item get(int64 index);
};
struct Numbers: Container
{
    type Item = int64;
    Item get(int64 index)
    {
        return index + 10;
    }
};
C.Item first<C: Container>(C value)
{
    return value.get(0);
}
int64 main()
{
    Numbers numbers;
    println("{}", first<Numbers>(numbers));
    return 0;
}
\end{lstlisting}
\par
\medskip

\par
\medskip
\begin{lstlisting}[style=out]
10
\end{lstlisting}
\par
\medskip

接口内是关联类型产生式，模块层的 \texttt{type} 仍是类型别名，原有类型参数 \texttt{<\allowbreak{}T:\allowbreak{} Valued>\allowbreak{}} 保持可用。

\subsection{泛型返回控制表达式}

返回块、\texttt{if}、\texttt{match} 和 \texttt{loop} 同样可以使用类型参数 T。具体类型确定后，编译器再按该类型检查每条结果路径。

\par
\medskip
\begin{lstlisting}[style=magel]
T fromLocal<T>(T value)
{
    return
    {
        T local = value;
        yield local;
    };
}

T chooseValue<T>(bool first, T left, T right)
{
    return if (first)
    {
        T local = left;
        yield local;
    }
    else
    {
        yield right;
    };
}

int64 main()
{
    String word = "generic";
    String copied = fromLocal<String>(word);
    int64 number = chooseValue<int64>(false, 4, 9);
    println("{} {} {}", word, copied, number);
    return 0;
}
\end{lstlisting}
\par
\medskip

\par
\medskip
\begin{lstlisting}[style=out]
generic generic 9
\end{lstlisting}
\par
\medskip

\texttt{local} 在值表达式内部，通过 \texttt{break} 隐式移动；\texttt{right} 在表达式外部，按值复制。函数形参本身也是调用者值的副本，所以这里不会消耗调用者的 \texttt{word}。泛型不会取消分支类型一致、生命周期或清理检查，返回引用也不能指向已经退出的局部对象。

\section{小结}

\begin{center}
\begin{tabular}{ll}
\toprule
你会了 & 关键点 \\
\midrule
泛型函数 & \texttt{T max\allowbreak{}Of<\allowbreak{}T>\allowbreak{}(\allowbreak{}T a,\allowbreak{} T b)\allowbreak{}}，类型参数写在名字后 \\
调用 & 可推导 \texttt{max\allowbreak{}Of(\allowbreak{}3,\allowbreak{} 9)\allowbreak{}}，可显式 \texttt{max\allowbreak{}Of<\allowbreak{}int64>\allowbreak{}(\allowbreak{}3,\allowbreak{} 9)\allowbreak{}} \\
泛型结构体 & \texttt{struct Duo<\allowbreak{}T,\allowbreak{} U>\allowbreak{} \{ T first; U second; \}} \\
约束 & \texttt{<\allowbreak{}T:\allowbreak{} Valued>\allowbreak{}}，保证函数体能调接口方法 \\
单态化 & 编译期按类型实参各生成一份，运行期无分派开销 \\
报错时机 & 约束不满足 \ensuremath{\rightarrow} 调用处报；无约束泛型的体 \ensuremath{\rightarrow} 实例化时报 \\
命名坑 & 别用标准库已占的泛型名（如 \texttt{Pair}） \\
默认类型参数 & \texttt{<\allowbreak{}T,\allowbreak{} U =\allowbreak{} T>\allowbreak{}}（进阶，先了解） \\
\bottomrule
\end{tabular}
\end{center}

\section{下一步}

到现在为止，你写的所有函数都得先有名字才能用：\texttt{max\allowbreak{}Of}、\texttt{read\_value}……可有时候你只想临时写一小段逻辑，随手传给别的函数——排序的「怎么比」、遍历的「对每个元素做什么」。

这就要把函数本身当值来传。第 15 章讲函数值、\texttt{fn} 和闭包。

翻到\hyperref[ch:15]{第 15 章 函数值与闭包}。
